\chapter{Supplementary Simulation: Shared-Protein Dependence}

This chapter supplements the dependence discussion in Chapter~\ref{chap:protein_results} with a graph-aware simulation study. It is not meant to replace the main empirical analysis of the protein--protein interaction data. Rather, it isolates the effect of shared-protein dependence in a controlled setting and shows how that dependence changes uncertainty quantification and labeled-sample calculations when the unit of analysis is an interaction edge. The point of the simulation is that the estimand and the prediction-powered workflow are held fixed while the strength of endpoint dependence can be varied directly.

\section{Simulation Design}

I simulate an undirected protein--protein interaction network with edge labels. Each Monte Carlo replicate begins with $m=320$ proteins. Protein $j$ receives a latent node effect
\[
a_j \sim N(0, \tau^2),
\]
where $\tau$ controls the strength of shared-protein dependence. The parameter is a standard deviation on the latent-score scale, so its admissible range is $\tau \in [0,\infty)$. In this simulation, $\tau=0$ corresponds to the independent-edge baseline, while larger values introduce stronger shared-endpoint dependence through the common node effects. I then sample $M=20{,}000$ unordered candidate edges $e=(u,v)$ from the complete graph on the 320 proteins. For each sampled edge, an independent edge-specific signal $z_e \sim N(0,1)$ is drawn, and the true logit signal is defined as
\[
\eta_e = \beta_0 + \beta_1 z_e + a_u + a_v,
\]
with $(\beta_0,\beta_1)=(-2.2,1.0)$. The true binary edge label is then generated by
\[
Y_e \sim \mathrm{Bernoulli}\!\left(\operatorname{logit}^{-1}(\eta_e)\right).
\]
To build a noisy surrogate, I perturb the same latent signal with additional score noise,
\[
\hat p_e = \operatorname{logit}^{-1}(\eta_e + \varepsilon_e),
\qquad
\varepsilon_e \sim N(0, \sigma_{\mathrm{score}}^2),
\]
with $\sigma_{\mathrm{score}}=1.0$. This construction keeps the surrogate informative without making it deterministically tied to the truth.

The interpretation of $\tau$ is clearest on the latent scale. Consider two distinct edges that share one protein, such as $(u,v)$ and $(u,w)$. Because the edge-specific signals and the node effects for $v$ and $w$ are independent, the only common latent component is $a_u$. Under the simulation model,
\[
\Var(\eta_{uv}) = \beta_1^2 + 2\tau^2
\qquad\text{and}\qquad
\Cov(\eta_{uv}, \eta_{uw}) = \Var(a_u) = \tau^2.
\]
Hence the latent-score correlation for two edges sharing exactly one protein is
\[
\Corr(\eta_{uv}, \eta_{uw})
=
\frac{\tau^2}{\beta_1^2 + 2\tau^2}.
\]
With $\beta_1=1$, the three simulation settings $\tau \in \{0,0.5,1.0\}$ correspond to latent shared-endpoint correlations of approximately $0$, $0.17$, and $0.33$. The Bernoulli draw and logistic transform make the corresponding label dependence nonlinear, but the latent correlation still provides a useful scale for interpreting the strength of the induced dependence.

\begin{table}[ht]
\centering
\small
\begin{tabular}{@{}ll@{}}
\toprule
Setting & Value \\
\midrule
Number of proteins & 320 \\
Candidate edges per replicate & 20,000 \\
Prediction-only edges & 4,000 \\
Labeled-sample grid & 50 to 500 by 50; 750 to 2,000 by 250; 3,000 to 8,000 by 500; 9,000 to 16,000 by 1,000 \\
Monte Carlo replicates & 500 \\
Node-effect standard deviations & 0.00, 0.50, 1.00 \\
Effect-size grid & 0.02, 0.03, 0.04, 0.05, 0.06, 0.07, 0.08, 0.09, 0.10 \\
Logit intercept beta0 & -2.2 \\
Edge-signal coefficient beta1 & 1.0 \\
Surrogate-noise standard deviation & 1.0 \\
\bottomrule
\end{tabular}
\caption[Simulation settings for the endpoint-clustered study]{Simulation settings for the stand-alone endpoint-clustered dependence study.}
\label{tab:sim_endpoint_cluster_settings}
\end{table}


Within each replicate, I split the sampled edges into a prediction-only pool of size $N=4000$ and a nested labeled pool whose size ranges from 50 to 16,000 edges, with finer spacing at small budgets and coarser spacing at larger ones. The same prediction-only pool is used for every candidate $n$ within the replicate, while the labeled pool is revealed gradually. This design makes the planning curves comparable across a much wider range of labeled budgets. The full simulation settings, including the dependence levels and the number of Monte Carlo replicates, are listed in Table~\ref{tab:sim_endpoint_cluster_settings}.

The estimand is the edge-level mean
\[
\theta = \mathbb{E}[Y_e],
\]
which may be read as the interaction prevalence in the simulated target edge population. As in Chapters~\ref{chap:foundation_power} and~\ref{chap:protein_results}, I use \(f_e\) for the generic surrogate and then specialize to the protein-style score \(f_e=\hat p_e\). I compare two estimators: the labeled-only sample mean and the mean \texttt{PPI++} estimator. Vanilla prediction-powered inference corresponds to the special case \(\lambda=1\), but the simulation focuses on the tuned \texttt{PPI++} version because that is the main planning target in the earlier chapters. To keep the simulation focused on dependence rather than on tuning noise, the \texttt{PPI++} weight is taken to be the oracle weight implied by the simulated moments,
\[
\lambda^{\star}(n)
=
\frac{\Cov(Y_e, f_e)}{(1+n/N)\Var(f_e)},
\]
which is the same oracle choice used in the one-sample mean planning formulas reviewed in Chapter~\ref{chap:foundation_power} \parencite{ppi_pp, ppi_power}. In the simulation itself, \(f_e=\hat p_e\).

\section{Edge-Level Contributions and Endpoint-Clustered Variance}

The starting point is the same one-sample mean estimator family used earlier in the thesis. For the labeled-only mean,
\[
\hat\theta_{\mathrm{mean}}
=
\frac{1}{n}\sum_{e \in L} Y_e.
\]
For a fixed tuning parameter \(\lambda\), the mean \texttt{PPI++} estimator is
\[
\hat\theta_{\lambda}
=
\frac{1}{n}\sum_{e \in L}\{Y_e-\lambda f_e\}
+
\frac{\lambda}{N}\sum_{e \in U} f_e,
\]
and the simulation uses the oracle choice \(\lambda=\lambda^\star(n)\). Under the independent-edge working model from Chapter~\ref{chap:foundation_power}, the exact variance baselines are
\[
\Var_{\mathrm{iid}}(\hat\theta_{\mathrm{mean}})
=
\frac{\Var(Y_e)}{n}
\]
and
\[
\Var_{\mathrm{iid}}(\hat\theta_{\lambda})
=
\frac{\Var(Y_e-\lambda f_e)}{n}
+
\frac{\lambda^2 \Var(f_e)}{N}.
\]
These are the same split-sample formulas used in the earlier planning chapter, now specialized to the edge-level setting.

To construct a graph-aware analogue, it is convenient to rewrite the estimator as an average of edge-level contributions. If \(E_{\mathrm{obs}}=L\cup U\) denotes the combined observed edge set, then
\[
\hat\theta_{\lambda} = \frac{1}{|E_{\mathrm{obs}}|}\sum_{e\in E_{\mathrm{obs}}} X_e,
\]
where
\[
X_e =
\begin{cases}
\dfrac{n+N}{n}\{Y_e-\lambda f_e\}, & e \in L,\\[6pt]
\dfrac{n+N}{N}\lambda f_e, & e \in U.
\end{cases}
\]
This representation is useful because it rewrites the split-sample estimator as a single average without changing the estimator itself.

Shared-protein dependence alters the variance calculation because two edges can move together when they share a protein. To reflect that structure, let \(\tilde X_e = X_e - \hat\theta_{\lambda}\) and define the endpoint-cluster sum for protein \(j\) by
\[
C_j = \sum_{e\ni j} \tilde X_e,
\]
where the sum runs over all observed edges incident to protein $j$. I then use the endpoint-clustered variance estimate
\[
\widehat{\Var}_{\mathrm{end}}(\hat\theta_{\lambda})
=
\frac{1}{|E_{\mathrm{obs}}|^2}
\left(
\sum_{j=1}^{m} C_j^2 - \sum_{e\in E_{\mathrm{obs}}} \tilde X_e^2
\right).
\]
This is the scalar, intercept-only analogue of the inclusion--exclusion logic used in multiway clustered sandwich estimators \parencite{cameron2011multiway}, adapted to the dyadic setting in which each observed edge belongs to exactly two endpoint clusters. It is also closely related to the dyadic cluster-robust variance viewpoint developed by \textcite{aronow2015dyadic}, who emphasize that dyads sharing a member can have correlated errors and therefore require cluster-robust variance adjustment. The subtraction term is needed because
\[
\sum_{j=1}^{m} C_j^2
=
\sum_{j=1}^{m}\sum_{e\ni j}\sum_{e'\ni j} \tilde X_e \tilde X_{e'}.
\]
In this expansion, each edge contributes its own squared term twice, once for each of its two endpoints, while a distinct pair of edges contributes a cross-product once if and only if the two edges share an endpoint. Therefore
\[
\sum_{j=1}^{m} C_j^2 - \sum_{e\in E_{\mathrm{obs}}} \tilde X_e^2
\]
keeps each diagonal contribution once and adds the cross-products associated with shared endpoints. I do not present this as a new asymptotic theorem. Rather, it is a transparent graph-aware variance construction for the simulation setting, chosen to stay close to standard clustered-sandwich ideas while tracking the dependence induced by shared endpoints. In the reported summaries, the independent-edge baseline is always the exact split-sample variance above, and the endpoint-clustered expression is used as its graph-aware counterpart.

\section{Simulation Workflow}

The Monte Carlo experiment proceeds in five fixed steps.
\begin{enumerate}
  \item Draw protein-level node effects and sample a set of unordered edges.
  \item Generate the true edge labels $Y_e$ and the noisy surrogate scores $\hat p_e$.
  \item Split the sampled edges into a prediction-only pool and a nested labeled pool.
  \item For each labeled budget $n$, compute the labeled-only mean and mean \texttt{PPI++} estimators, together with their i.i.d. and endpoint-clustered variance estimates.
  \item Aggregate the replicate-level output into variance-inflation and planning summaries.
\end{enumerate}
The corresponding code is written in the stand-alone R driver \texttt{07\_simulate\_endpoint\_clustered\_planning.R} together with the new helper file \texttt{endpoint\_clustered\_utils.R}. No existing thesis scripts are modified.

\section{Monte Carlo Results on Variance Inflation}

\begin{figure}[ht]
\centering
\includegraphics[width=0.92\textwidth]{sim_endpoint_cluster_variance.png}
\caption[Simulation-based variance inflation under shared-protein dependence]{Monte Carlo comparison of endpoint-clustered and i.i.d. uncertainty estimates. The vertical axis reports the ratio of the endpoint-clustered standard error to the corresponding i.i.d. standard error. When $\tau=0$, the two estimates are nearly the same. As shared-protein dependence strengthens, the endpoint-clustered standard error increases above the i.i.d. baseline.}
\label{fig:sim_endpoint_cluster_variance}
\end{figure}

Figure~\ref{fig:sim_endpoint_cluster_variance} shows the main dependence pattern clearly. When $\tau=0$, the simulated edges are close to the independent-edge baseline, so the endpoint-clustered and i.i.d. standard errors nearly coincide. As $\tau$ increases, incident edges inherit more common signal through the node effects, and the endpoint-clustered standard error moves upward. This happens for both estimators, although the amount of inflation differs because the \texttt{PPI++} contribution mixes labeled residual-correction terms with prediction-only surrogate terms.

\begin{table}[ht]
\centering
\small
\begin{tabularx}{\textwidth}{@{}lccccccc@{}}
\toprule
Estimator & $\tau$ & Prevalence & Pilot $\rho^2$ & $\lambda^{\star}$ & i.i.d. SE & Endpoint SE & SE ratio \\
\midrule
Labeled mean & 0.00 & 0.134 & 0.052 & -- & 0.024 & 0.024 & 0.990 \\
PPI++ & 0.00 & 0.134 & 0.052 & 0.425 & 0.024 & 0.024 & 1.005 \\
Labeled mean & 0.50 & 0.148 & 0.093 & -- & 0.025 & 0.026 & 1.020 \\
PPI++ & 0.50 & 0.148 & 0.093 & 0.529 & 0.024 & 0.025 & 1.032 \\
Labeled mean & 1.00 & 0.185 & 0.202 & -- & 0.027 & 0.030 & 1.082 \\
PPI++ & 1.00 & 0.185 & 0.202 & 0.695 & 0.025 & 0.027 & 1.110 \\
\bottomrule
\end{tabularx}
\caption[Summary of simulated endpoint-clustered results at n=200]{Summary of the simulation output at labeled sample size $n=200$. The SE ratio is the endpoint-clustered standard error divided by the corresponding i.i.d. standard error.}
\label{tab:sim_endpoint_cluster_summary}
\end{table}


Table~\ref{tab:sim_endpoint_cluster_summary} gives one concrete slice of the simulation at labeled sample size $n=200$. The table reports the simulated prevalence, the pilot agreement measure \(\rho_{Y\hat p}^2=\Corr(Y_e,\hat p_e)^2\), which is the protein-specific version of the generic \(\rho_{Yf}^2\), the oracle weight \(\lambda^{\star}\) for \texttt{PPI++}, and the resulting i.i.d. and endpoint-clustered standard errors. The final column is the key comparison. It is close to 1 when the node effect is absent, and it increases as shared-protein dependence becomes stronger.

\section{Implications for Labeled-Sample Calculations}

To connect the simulation output to planning, I use the same variance-threshold logic as in the main planning discussion. For a target absolute difference $\Delta$, significance level $\alpha$, and power $1-\beta$, define
\[
S^2
=
\left(
\frac{\Delta}{z_{1-\alpha/2}+z_{1-\beta}}
\right)^2.
\]
Rather than showing only the thresholded sample-size output, it is more informative to plot the corresponding minimum detectable absolute difference as a function of $n$:
\[
\Delta_{\min}(n)
=
\bigl(z_{1-\alpha/2}+z_{1-\beta}\bigr)\sqrt{\widehat V(\hat\theta_n)}.
\]
Here \(\widehat V(\hat\theta_n)\) denotes either the exact i.i.d.\ plug-in variance or the endpoint-clustered variance estimate for the estimator being plotted at labeled budget \(n\). This uses the entire variance curve rather than only the first grid point that crosses a threshold. It therefore shows more clearly how dependence changes planning over a wide range of labeled budgets.

\begin{figure}[ht]
\centering
\includegraphics[width=0.96\textwidth]{sim_endpoint_cluster_mde.png}
\caption[Simulation-based minimum detectable effects under endpoint dependence]{Minimum detectable absolute difference as a function of labeled sample size for the i.i.d. and endpoint-clustered variance rules. The horizontal axis is shown on a log scale to display the wide labeled-sample range clearly. As shared-protein dependence strengthens, the endpoint-clustered rule yields larger detectable-effect thresholds, especially for the stronger-dependence settings and the prediction-powered estimator.}
\label{fig:sim_endpoint_cluster_mde}
\end{figure}

Figure~\ref{fig:sim_endpoint_cluster_mde} turns the variance comparison into a planning comparison without collapsing the result to a single thresholded $n$. Under weak or absent dependence, the i.i.d. and endpoint-clustered curves are close. Once the node effect is large enough to create visible overlap dependence, the endpoint-clustered curve remains above the i.i.d. curve across the labeled-budget range. The same general conclusion holds for both estimators: prediction-powered structure can still reduce the labeled burden, but a graph-aware variance rule attenuates the apparent gain when edges are no longer effectively independent.

\section{Takeaway}

This simulation gives a more concrete dependence story than a scalar variance-inflation knob alone. Because the dependence is induced by shared endpoints in an explicit edge-generating mechanism, the resulting uncertainty and planning summaries can be read directly against the graph structure. The main lesson is not that one variance rule is universally correct. Rather, once edge dependence is introduced through shared proteins, an edge-level mean estimator should be evaluated with a dependence-aware variance construction rather than with an independent-edge formula alone.
